\section{Preliminaries}
\label{sec:prelims}

We fix a finite set $\vars$ of program variables and a countable state space $\Sigma$; a state $\sigma\in\Sigma$ maps variables to rational values. Arithmetic and Boolean expressions are interpreted as functions on states. For a Boolean expression $b$, let $\llbracket b\rrbracket=\{\sigma\mid \sigma\models b\}$ and let $[b]$ denote its Iverson bracket. Updating $x$ in $\sigma$ to $a(\sigma)$ is written $\sigma[x/a(\sigma)]$.

\subsection{Probabilistic programs and subdistributions}

We consider the probabilistic guarded-command language
\[
 C ::= \ski \mid \ass{x}{a} \mid \comp{C}{C} \mid \pchoice{C}{p}{C}
      \mid \ifelse{b}{C}{C} \mid \while{b}{C},
\]
where $p\in[0,1]\cap\mathbb{Q}$. A \emph{subdistribution} is a function $\mu:\Sigma\to\mathbb{R}_{\geq0}$ of weight $\weight{\mu}=\sum_{\sigma}\mu(\sigma)\leq1$. We write $\subdists$ for all subdistributions, $\supp(\mu)=\{\sigma\mid\mu(\sigma)>0\}$ for the support, and $\dirac{\sigma}$ for the Dirac distribution at $\sigma$. Missing mass represents nontermination.

Programs denote linear post-distribution transformers $\post{C}$. The clauses needed below are
\begin{align*}
 \post{\ski}(\mu) &= \mu,\\
 \post{\ass{x}{a}}(\mu)(\sigma)
   &= \sum_{\sigma'=\sigma[x/a(\sigma')]}\mu(\sigma'),\\
 \post{\pchoice{C_1}{p}{C_2}}(\mu)
   &= \post{C_1}(p\mu)+\post{C_2}((1-p)\mu),\\
 \post{\comp{C_1}{C_2}}(\mu)
   &= \post{C_2}(\post{C_1}(\mu)),\\
 \post{\ifelse{b}{C_1}{C_2}}(\mu)
   &= \post{C_1}([b]\mu)+\post{C_2}([\neg b]\mu).
\end{align*}
For a loop $C=\while{b}{B}$, let $F_\mu(e)=\mu+\post{B}([b]e)$. Its semantics is
\[
  \post{C}(\mu)=[\neg b]\cdot \lfp F_\mu.
\]
Equivalently, the Kleene approximants collect executions that leave the loop after finitely many iterations. Linearity gives
$\post{C}(\sum_i p_i\dirac{\sigma_i})=\sum_i p_i\post{C}(\dirac{\sigma_i})$, which lets us reduce one-iteration obligations to Dirac inputs.

\begin{definition}[Almost-sure termination]
A program $C$ is \emph{almost-surely terminating} on $M\subseteq\subdists$ if
$\weight{\post{C}(\mu)}=\weight{\mu}$ for every $\mu\in M$.
\end{definition}

Thus AST means that execution loses no probability mass. We do not require positive AST: the expected runtime may be infinite.

\subsection{Distributional invariants}

Our assertions are sets of distributions rather than predicates over individual states. This aligns the termination argument with existing partial-correctness proofs for sampling algorithms~\cite{zilken2026verifyingsamplingalgorithmsdistributional}.

\begin{definition}[Distributional invariant]
Let $C=\while{b}{B}$ and $M\subseteq\subdists$. A set $I\subseteq\subdists$ is a \emph{distributional invariant} for $C$ and $M$ if
\[
 M\subseteq I
 \quad\text{and}\quad
 [\neg b]\mu+\post{B}([b]\mu)\in I
 \quad\text{for all }\mu\in I.
\]
We abbreviate $\supp(I)=\bigcup_{\mu\in I}\supp(\mu)$.
\end{definition}

The invariant can simultaneously carry state-space bounds and genuinely distributional information such as equal probability of states in a row of a decision tree. The AST rule uses the support information, while retaining the stronger invariant permits direct composition with partial correctness.

\subsection{Supermartingales, variants, and sublevel sets}

A function $v:\supp(I)\to\mathbb{R}_{\geq0}$ is a one-iteration supermartingale for loop body $B$ if, for every enabled state $\sigma$,
\[
  \sum_{\sigma'\in\Sigma}\post{B}(\dirac{\sigma})(\sigma')v(\sigma')\leq v(\sigma).
\]
For $r\in\mathbb{R}_{\geq0}$, its sublevel set is
$v_{\leq r}=\{\sigma\in\supp(I)\mid v(\sigma)\leq r\}$.
A variant $u:\supp(I)\to\mathbb{N}$ need not decrease in expectation or on every outcome. We only require a uniformly positive chance to move to a smaller value. The role of $v$ is to ensure that executions cannot escape all regions in which $u$ has a finite upper bound.

This division of labour matters for rejection samplers. A rejected run may reset $u$ to a larger value; AST follows because progress remains possible with nonzero probability while $v$ controls the size of the relevant search space.
